For $\Qx$, the norm plays the role that size plays in integer division.  Since norms of nonzero
polynomials are positive integers, we can use the well-ordering principle to
pick a remainder of smallest norm.

\begin{theorem}[Division algorithm in $\Qx$]
Let $f(x),g(x)\in\Qx$ with $g(x)\neq0$.  Then there exist unique
polynomials $q(x),r(x)\in\Qx$ such that
\[
f(x)=q(x)g(x)+r(x),
\]
where either
\[
r(x)=0
\]
or
\[
\|r(x)\|<\|g(x)\|.
\]
\end{theorem}

\begin{proof}
\textbf{Existence.}
Look at all possible remainders after subtracting multiples of $g$ from $f$:
\[
S=\{\,f(x)-q(x)g(x):q(x)\in\Qx\,\}.
\]
If $0\in S$, then $f=qg$ for some $q$, so
\[
f=qg+0,
\]
and the result holds with $r=0$.

Now suppose $0\notin S$.  Then every element of $S$ is nonzero, so
each element has a norm in $\mathbb{Z}^+$.  Define
\[
D=\{\,\|s(x)\|:s(x)\in S\,\}\subseteq\mathbb{Z}^+.
\]
The set $D$ is nonempty because $f-0g=f$ lies in $S$.  By the well-ordering
principle, $D$ has a smallest element.  Choose $r\in S$ with minimum norm.
Since $r\in S$, there exists $q\in\Qx$ such that
\[
r=f-qg.
\]
Equivalently,
\[
f=qg+r.
\]

It remains to show that $\|r(x)\|<\|g(x)\|$.  Suppose instead that
\[
\|r(x)\|\geq\|g(x)\|.
\]
Write
\[
r(x)=a_mx^m+\cdots
\qquad\text{and}\qquad
g(x)=b_nx^n+\cdots,
\]
where $a_m,b_n\neq0$, $m=\deg(r)$, and $n=\deg(g)$.  Since
$\|r(x)\|\geq\|g(x)\|$ means $2^m\geq2^n$, we get $m\geq n$.  Since $b_n\neq0$
and $b_n\in\mathbb{Q}$, the rational number $a_m/b_n$ exists.  Define
\[
r^*(x)=r(x)-\frac{a_m}{b_n}x^{m-n}g(x).
\]
The second term has leading term $a_mx^m$, so the leading term of $r$ cancels.
Thus either $r^*=0$ or
\[
\|r^*(x)\|<\|r(x)\|.
\]
Also,
\[
\begin{aligned}
r^*
&=f-qg-\frac{a_m}{b_n}x^{m-n}g\\
&=f-\left(q+\frac{a_m}{b_n}x^{m-n}\right)g.
\end{aligned}
\]
Therefore $r^*\in S$.  If $r^*=0$, then $0\in S$, which is not the case we are
in.  If $r^*\neq0$, then $r^*$ is in $S$ and has smaller norm than $r$,
contradicting how we chose $r$.  Therefore our assumption was false, and
\[
\|r(x)\|<\|g(x)\|.
\]

\textbf{Uniqueness.}
Suppose
\[
f=q_1g+r_1=q_2g+r_2,
\]
where $r_1$ and $r_2$ are both either zero or have norm less than $\|g(x)\|$.
Subtracting gives
\[
g(q_1-q_2)=r_2-r_1.
\]
Assume $q_1\neq q_2$.  Then $q_1-q_2\neq0$, and since $\Qx$ has no
zero divisors,
\[
g(q_1-q_2)\neq0.
\]
Hence
\[
\begin{aligned}
\|g(x)(q_1(x)-q_2(x))\|
&=2^{\deg(g(q_1-q_2))}\\
&=2^{\deg(g)+\deg(q_1-q_2)}\\
&=\|g(x)\|\cdot\|q_1(x)-q_2(x)\|\\
&\geq\|g(x)\|.
\end{aligned}
\]
But $r_2-r_1$ is either zero or has norm strictly less than $\|g(x)\|$, since
each of $r_1$ and $r_2$ is either zero or has norm less than $\|g(x)\|$, so
their degrees are both less than $\deg(g)$ when they are nonzero.  This is
impossible.
Therefore $q_1=q_2$.  Substituting back gives $r_1=r_2$.
\end{proof}
%%TODO: GCD
